% !TeX root = ../../book.tex

%
% Custom theorem environments
%

% Styles

\declaretheoremstyle[
    spaceabove=6pt,
    spacebelow=6pt,
    headfont={\color{thmcol}\fontfamily{bch}\bfseries\hspace{-\thmsymbolwidth}\thmintrosymbol},
    headpunct={{\par\nobreak}\\},
    notefont={\mdseries},
    notebraces={(}{)},
    bodyfont=\normalfont
]{infdescthm}

\declaretheoremstyle[
    spaceabove=6pt,
    spacebelow=6pt,
    headfont={\color{thmcol}\fontfamily{bch}\bfseries\hspace{-\thmsymbolwidth}\thmintrosymbol},
    headpunct={{\par\nobreak}\\},
    notefont={\mdseries},
    notebraces={(}{)},
    bodyfont={\fixthmbox\normalfont},
    shaded={bgcolor=thmbgcol, rulecolor=thmbdcol, rulewidth=0.25pt}
]{infdescimptthm}

\declaretheoremstyle[
    spaceabove=6pt,
    spacebelow=6pt,
    headfont={\color{defcol}\fontfamily{bch}\bfseries\hspace{-\defsymbolwidth}\defintrosymbol},
    headpunct={{\par\nobreak}\\},
    notefont={\mdseries},
    notebraces={(}{)},
    bodyfont=\normalfont
]{infdescdef}

\declaretheoremstyle[
    spaceabove=6pt,
    spacebelow=6pt,
    headfont={\color{defcol}\fontfamily{bch}\bfseries\hspace{-\defsymbolwidth}\defintrosymbol},
    headpunct={{\par\nobreak}\\},
    notefont={\mdseries},
    notebraces={(}{)},
    bodyfont={\fixthmbox\normalfont},
    shaded={bgcolor=defbgcol, rulecolor=defbdcol, rulewidth=0.25pt}
]{infdescimptdef}

\declaretheoremstyle[
    spaceabove=6pt,
    spacebelow=6pt,
    headfont={\color{excol}\fontfamily{bch}\bfseries\hspace{-\exsymbolwidth}\exintrosymbol},
    headpunct={{\par\nobreak}\\},
    notefont={\mdseries},
    notebraces={(}{)},
    bodyfont=\normalfont,
    qed={\exqedsymbol}
]{infdescex}

\declaretheoremstyle[
    spaceabove=6pt,
    spacebelow=6pt,
    headfont={\color{excol}\fontfamily{bch}\bfseries\hspace{-\exsymbolwidth}\exintrosymbol},
    headpunct={{\par\nobreak}\\},
    notefont={\mdseries},
    notebraces={(}{)},
    bodyfont={\fixthmbox\normalfont},
    shaded={bgcolor=exbgcol, rulecolor=exbdcol, rulewidth=0.25pt}
]{infdescimptex}

\declaretheoremstyle[
    spaceabove=6pt,
    spacebelow=6pt,
    headfont={\color{excol}\fontfamily{bch}\bfseries\hspace{-\quotesymbolwidth}\quoteintrosymbol},
    headpunct={{\par\nobreak}\\},
    notefont={\mdseries},
    notebraces={(}{)},
    bodyfont=\normalfont,
    qed={\quoteqedsymbol}
]{infdescquote}

\declaretheoremstyle[
    spaceabove=6pt,
    spacebelow=6pt,
    headfont={\color{prcol}\fontfamily{bch}\bfseries\hspace{-\prsymbolwidth}\printrosymbol},
    headpunct={{\par\nobreak}\\},
    notefont={\mdseries},
    notebraces={(}{)},
    bodyfont=\normalfont,
    qed={\exqedsymbol}
]{infdescpr}

\declaretheoremstyle[
    spaceabove=6pt,
    spacebelow=6pt,
    headfont={\color{prcol}\fontfamily{bch}\bfseries\hspace{-\prsymbolwidth}\printrosymbol},
    headpunct={{\par\nobreak}\\},
    notefont={\mdseries},
    notebraces={(}{)},
    bodyfont={\fixthmbox\normalfont},
    shaded={bgcolor=prbgcol, rulecolor=prbdcol, rulewidth=0.25pt}
]{infdescimptpr}

\declaretheoremstyle[
    spaceabove=6pt,
    spacebelow=6pt,
    headfont={\color{midgray}\fontfamily{bch}\bfseries},
    headpunct={.},
    notefont={\mdseries},
    notebraces={(}{)},
    bodyfont=\normalfont
]{infdescchapex}

\declaretheoremstyle[
    spaceabove=6pt,
    spacebelow=6pt,
    headfont={\color{tipcol}\fontfamily{bch}\bfseries\hspace{-\tipsymbolwidth}\tipintrosymbol},
    headpunct={\par\nobreak\\},
    notefont={\mdseries},
    notebraces={(}{)},
    bodyfont=\normalfont,
    qed={\tipqedsymbol}
]{infdesctip}

\declaretheoremstyle[
    spaceabove=6pt,
    spacebelow=6pt,
    headfont={\color{tipcol}\fontfamily{bch}\bfseries\hspace{-\tipsymbolwidth}\tipintrosymbol},
    headpunct={\par\nobreak\\},
    notefont={\mdseries},
    notebraces={(}{)},
    bodyfont={\fixthmbox\normalfont},
    shaded={bgcolor=tipbgcol, rulecolor=tipbdcol, rulewidth=0.25pt}
]{infdescimpttip}

\declaretheoremstyle[
    spaceabove=6pt,
    spacebelow=6pt,
    headfont={\color{pfcol}\fontfamily{bch}\bfseries\itshape},
    headpunct={\par\nobreak\\},
    notefont={\mdseries},
    notebraces={}{},
    bodyfont=\normalfont,
    qed={\color{pfcol}$\Box$}
]{infdescpf}

\declaretheoremstyle[
    spaceabove=6pt,
    spacebelow=6pt,
    headfont={\color{pfcol}\fontfamily{bch}\bfseries\itshape},
    headpunct={\par\nobreak\\},
    notefont={\mdseries},
    notebraces={}{},
    bodyfont=\normalfont
]{infdescpfnoqed}



% Results

\declaretheorem[style=infdescimptthm, numberwithin=section, Refname={定理}, name=定理]{theorem}
\declaretheorem[style=infdescimptthm, sibling=theorem, Refname={定理}, name={定理\assumesAC{thmcolac}}]{theoremac}
\declaretheorem[style=infdescimptthm, numberwithin=chapter, Refname={定理}, name=定理]{chtheorem}
\declaretheorem[style=infdescimptthm, sibling=theorem, refname={定理}, Refname={定理}, name=定理]{itheorem}
\declaretheorem[style=infdescimptthm, numbered=no, refname={定理}, Refname={定理}, name=定理]{theorem*}
\declaretheorem[style=infdescimptthm, sibling=theorem, refname={定理}, Refname={定理}, name={\optmark{定理}}]{otheorem}

\declaretheorem[style=infdescthm, sibling=theorem, refname={引理}, Refname={引理}, name=引理]{lemma}
\declaretheorem[style=infdescthm, sibling=theorem, refname={引理}, Refname={引理}, name={引理\assumesAC{thmcolac}}]{lemmaac}
\declaretheorem[style=infdescthm, sibling=chtheorem, refname={引理}, Refname={引理}, name=引理]{chlemma}
\declaretheorem[style=infdescimptthm, sibling=theorem, refname={引理}, Refname={引理}, name=引理]{ilemma}
\declaretheorem[style=infdescthm, numbered=no, refname={引理}, Refname={引理}, name=引理]{lemma*}
\declaretheorem[style=infdescthm, sibling=theorem, refname={引理}, Refname={引理}, name={\optmark{引理}}]{olemma}

\declaretheorem[style=infdescthm, sibling=theorem, refname={命题}, Refname={命题}, name=命题]{proposition}
\declaretheorem[style=infdescthm, sibling=theorem, refname={命题}, Refname={命题}, name={命题\assumesAC{thmcolac}}]{propositionac}
\declaretheorem[style=infdescthm, sibling=chtheorem, refname={命题}, Refname={命题}, name=命题]{chproposition}
\declaretheorem[style=infdescimptthm, sibling=theorem, refname={命题}, Refname={命题}, name=命题]{iproposition}
\declaretheorem[style=infdescthm, numbered=no, refname={命题}, Refname={命题}, name=命题]{proposition*}
\declaretheorem[style=infdescthm, sibling=theorem, refname={命题}, Refname={命题}, name={\optmark{命题}}]{oproposition}

\declaretheorem[style=infdescthm, sibling=theorem, refname={推论}, Refname={推论}, name=推论]{corollary}
\declaretheorem[style=infdescthm, sibling=theorem, refname={推论}, Refname={推论}, name={推论\assumesAC{thmcolac}}]{corollaryac}
\declaretheorem[style=infdescthm, sibling=chtheorem, refname={推论}, Refname={推论}, name=推论]{chcorollary}
\declaretheorem[style=infdescimptthm, sibling=theorem, refname={推论}, Refname={推论}, name=推论]{icorollary}
\declaretheorem[style=infdescthm, numbered=no, refname={推论}, Refname={推论}, name=推论]{corollary*}
\declaretheorem[style=infdescthm, sibling=theorem, refname={推论}, Refname={推论}, name={\optmark{推论}}]{ocorollary}

\declaretheorem[style=infdescimptthm, sibling=theorem, refname={公理}, Refname={公理}, name=公理]{axiom}
\declaretheorem[style=infdescimptthm, sibling=chtheorem, refname={公理}, Refname={公理}, name=公理]{chaxiom}
\declaretheorem[style=infdescimptthm, sibling=theorem, refname={公理}, Refname={公理}, name=公理]{iaxiom}
\declaretheorem[style=infdescimptthm, numbered=no, refname={公理}, Refname={公理}, name=公理]{axiom*}
\declaretheorem[style=infdescimptthm, sibling=theorem, refname={公理}, Refname={公理}, name={\optmark{公理}}]{oaxiom}

\declaretheorem[style=infdescimptthm, sibling=theorem, refname={公理}, Refname={公理}, name=公理]{axioms}
\declaretheorem[style=infdescimptthm, sibling=chtheorem, refname={公理}, Refname={公理}, name=公理]{chaxioms}
\declaretheorem[style=infdescimptthm, sibling=theorem, refname={公理}, Refname={公理}, name=公理]{iaxioms}
\declaretheorem[style=infdescimptthm, numbered=no, refname={公理}, Refname={公理}, name=公理]{axioms*}
\declaretheorem[style=infdescimptthm, sibling=theorem, refname={公理}, Refname={公理}, name={\optmark{公理}}]{oaxioms}



% Definitions

\declaretheorem[style=infdescimptdef, sibling=theorem, refname={定义}, Refname={定义}, name={定义}]{definition}
\declaretheorem[style=infdescimptdef, sibling=theorem, refname={定义}, Refname={定义}, name={定义\assumesAC{defcolac}}]{definitionac}
\declaretheorem[style=infdescimptdef, sibling=chtheorem, refname={定义}, Refname={定义}, name={定义}]{chdefinition}
\declaretheorem[style=infdescimptdef, sibling=theorem, refname={定义}, Refname={定义}, name=定义]{idefinition}
\declaretheorem[style=infdescimptdef, numbered=no, refname={定义}, Refname={定义}, name=定义]{definition*}
\declaretheorem[style=infdescimptdef, sibling=theorem, refname={定义}, Refname={定义}, name={\optmark{定义}}]{odefinition}

\declaretheorem[style=infdescimptdef, sibling=theorem, refname={构造}, Refname={构造}, name=构造]{construction}
\declaretheorem[style=infdescimptdef, sibling=theorem, refname={构造}, Refname={构造}, name={构造\assumesAC{defcolac}}]{constructionac}
\declaretheorem[style=infdescimptdef, sibling=chtheorem, refname={构造}, Refname={构造}, name=构造]{chconstruction}
\declaretheorem[style=infdescimptdef, sibling=theorem, refname={构造}, Refname={构造}, name=构造]{iconstruction}
\declaretheorem[style=infdescimptdef, numbered=no, refname={构造}, Refname={构造}, name=构造]{construction*}
\declaretheorem[style=infdescimptdef, sibling=theorem, refname={构造}, Refname={构造}, name={\optmark{构造}}]{oconstruction}

\declaretheorem[style=infdescimptdef, sibling=theorem, refname={定义}, Refname={定义}, name={定义}]{predefinition}
\declaretheorem[style=infdescimptdef, sibling=chtheorem, refname={定义}, Refname={定义}, name={定义}]{chpredefinition}
\declaretheorem[style=infdescimptdef, sibling=theorem, refname={定义}, Refname={定义}, name={定义}]{ipredefinition}
\declaretheorem[style=infdescimptdef, numbered=no, refname={定义}, Refname={定义}, name={定义}]{predefinition*}
\declaretheorem[style=infdescimptdef, sibling=theorem, refname={定义}, Refname={定义}, name={\optmark{定义}}]{opredefinition}

\declaretheorem[style=infdescdef, sibling=theorem, refname={注释}, Refname={注释}, name=注释]{notation}
\declaretheorem[style=infdescdef, sibling=chtheorem, refname={注释}, Refname={注释}, name=注释]{chnotation}
\declaretheorem[style=infdescimptdef, sibling=theorem, refname={注释}, Refname={注释}, name=注释]{inotation}
\declaretheorem[style=infdescdef, numbered=no, refname={注释}, Refname={注释}, name=注释]{notation*}
\declaretheorem[style=infdescdef, sibling=theorem, refname={注释}, Refname={注释}, name={\optmark{注释}}]{onotation}

\declaretheorem[style=infdescimptdef, sibling=theorem, refname={词汇}, Refname={词汇}, name={词汇}]{vocabulary}
\declaretheorem[style=infdescimptdef, sibling=chtheorem, refname={词汇}, Refname={词汇}, name={词汇}]{chvocabulary}
\declaretheorem[style=infdescimptdef, sibling=theorem, refname={词汇}, Refname={词汇}, name=词汇]{ivocabulary}
\declaretheorem[style=infdescimptdef, numbered=no, refname={词汇}, Refname={词汇}, name=词汇]{vocabulary*}
\declaretheorem[style=infdescimptdef, sibling=theorem, refname={词汇}, Refname={词汇}, name={\optmark{词汇}}]{ovocabulary}

\declaretheorem[style=infdescdef, sibling=theorem, refname={假设}, Refname={假设}, name=假设]{assumption}
\declaretheorem[style=infdescdef, sibling=chtheorem, refname={假设}, Refname={假设}, name=假设]{chassumption}
\declaretheorem[style=infdescimptdef, sibling=theorem, refname={假设}, Refname={假设}, name=假设]{iassumption}
\declaretheorem[style=infdescdef, numbered=no, refname={假设}, Refname={假设}, name=假设]{assumption*}
\declaretheorem[style=infdescdef, sibling=theorem, refname={假设}, Refname={假设}, name={\optmark{假设}}]{oassumption}



% Exercises and examples

\declaretheorem[style=infdescex, sibling=theorem, refname={示例}, Refname={示例}, name=示例]{example}
\declaretheorem[style=infdescex, sibling=theorem, refname={示例}, Refname={示例}, name={示例\assumesAC{excolac}}]{exampleac}
\declaretheorem[style=infdescex, sibling=chtheorem, refname={示例}, Refname={示例}, name=示例]{chexample}
\declaretheorem[style=infdescimptex, sibling=theorem, refname={示例}, Refname={示例}, name=示例]{iexample}
\declaretheorem[style=infdescex, numbered=no, refname={示例}, Refname={示例}, name=示例]{example*}
\declaretheorem[style=infdescex, sibling=theorem, refname={示例}, Refname={示例}, name={\optmark{示例}}]{oexample}

\declaretheorem[style=infdescpr, sibling=theorem, refname={练习}, Refname={练习}, name={练习}]{exercise}
\declaretheorem[style=infdescpr, sibling=theorem, refname={练习}, Refname={练习}, name={Exercise\assumesAC{excolac}}]{exerciseac}
\declaretheorem[style=infdescpr, sibling=chtheorem, refname={练习}, Refname={练习}, name={练习}]{chexercise}
\declaretheorem[style=infdescimptpr, sibling=theorem, refname={练习}, Refname={练习}, name=练习]{iexercise}
\declaretheorem[style=infdescpr, numbered=no, refname={练习}, Refname={练习}, name=练习]{exercise*}
\declaretheorem[style=infdescpr, sibling=theorem, refname={练习}, Refname={练习}, name={\optmark{练习}}]{oexercise}

\declaretheorem[style=infdescpr, sibling=theorem, refname={讨论}, Refname={讨论}, name=讨论]{discussion}
\declaretheorem[style=infdescpr, sibling=chtheorem, refname={讨论}, Refname={讨论}, name=讨论]{chdiscussion}
\declaretheorem[style=infdescimptpr, sibling=theorem, refname={讨论}, Refname={讨论}, name=讨论]{idiscussion}
\declaretheorem[style=infdescpr, numbered=no, refname={讨论}, Refname={讨论}, name=讨论]{discussion*}
\declaretheorem[style=infdescpr, sibling=theorem, refname={讨论}, Refname={讨论}, name={\optmark{讨论}}]{odiscussion}

\declaretheorem[style=infdescpr, sibling=theorem, refname={问题}, Refname={问题}, name=问题]{problem}
\declaretheorem[style=infdescpr, sibling=chtheorem, refname={问题}, Refname={问题}, name=问题]{chproblem}
\declaretheorem[style=infdescimptpr, sibling=theorem, refname={问题}, Refname={问题}, name=问题]{iproblem}
\declaretheorem[style=infdescpr, numbered=no, refname={问题}, Refname={问题}, name=问题]{problem*}
\declaretheorem[style=infdescpr, sibling=theorem, refname={问题}, Refname={问题}, name={\optmark{问题}}]{oproblem}

\declaretheorem[style=infdescquote, sibling=theorem, refname={摘录}, Refname={摘录}, name=摘录]{extract}
\declaretheorem[style=infdescquote, sibling=theorem, refname={摘录}, Refname={摘录}, name={Extract\assumesAC{excolac}}]{extractac}
\declaretheorem[style=infdescquote, sibling=chtheorem, refname={摘录}, Refname={摘录}, name=摘录]{chextract}
\declaretheorem[style=infdescquote, numbered=no, refname={摘录}, Refname={摘录}, name=摘录]{extract*}
\declaretheorem[style=infdescquote, sibling=theorem, refname={摘录}, Refname={摘录}, name={\optmark{摘录}}]{oextract}

% End-of-chapter exercises

\declaretheorem[style=infdescchapex, numberwithin=chapter, refname={问题}, Refname={问题}, name={}]{chapex}

% Command for referring to exercises in other chapters
\newcommand{\CrefCQ}[2]{%
    \Cref*{#1} \Cref{#2}%
}

\newcommand{\CrefTO}[2]{%
    \Cref{#1} 至 \Cref{#2}%
}

% Tips and special remarks

\declaretheorem[style=infdesctip, numbered=no, name=解题提示]{problemtip}
\declaretheorem[style=infdesctip, numbered=no, name={证明提示}]{prooftip}
\declaretheorem[style=infdesctip, numbered=no, name={书写提示}]{writingtip}
\declaretheorem[style=infdesctip, numbered=no, name={\LaTeX{} 提示}]{latextip}
\declaretheorem[style=infdesctip, numbered=no, name=格式提示]{formattip}
\declaretheorem[style=infdesctip, numbered=no, name=目标]{goal}
\declaretheorem[style=infdesctip, numbered=no, name=旁白]{aside}
\declaretheorem[style=infdesctip, numbered=no, name=常见错误]{commonerror}

\declaretheorem[style=infdesctip, sibling=theorem, refname={惯例}, Refname={惯例}, name=惯例]{convention}
\declaretheorem[style=infdesctip, sibling=chtheorem, refname={惯例}, Refname={惯例}, name=惯例]{chconvention}

\declaretheorem[style=infdesctip, sibling=theorem, refname={备注}, Refname={备注}, name=备注]{remark}
\declaretheorem[style=infdesctip, sibling=chtheorem, refname={备注}, Refname={备注}, name=备注]{chremark}

% Proofs and strategies

\declaretheorem[style=infdescpf, numbered=no, name=证明]{cproof}
\declaretheorem[style=infdescpf, numbered=no, name=证明草图]{csketch}
\declaretheorem[style=infdescpfnoqed, numbered=no, name=证明思路]{cidea}

\declaretheorem[style=infdescimpttip, sibling=theorem, refname={策略}, Refname={策略}, name={策略}]{strategy}
\declaretheorem[style=infdescimpttip, sibling=theorem, refname={策略}, Refname={策略}, name={策略\assumesAC{tipcolac}}]{strategyac}
\declaretheorem[style=infdescimpttip, sibling=chtheorem, refname={策略}, Refname={策略}, name={策略}]{chstrategy}
\declaretheorem[style=infdesctip, numbered=no, refname={策略}, Refname={策略}, name={策略}]{strategy*}

\declaretheorem[style=infdescimpttip, refname={deadly sin,deadly sins}, Refname={Deadly Sin, Deadly Sins}, name={Deadly Sin}]{deadlysin}
\renewcommand*{\thedeadlysin}{the \Ordinalstring{deadlysin}}

\declaretheorem[style=infdescimpttip, sibling=theorem, refname={书写原则}, Refname={书写原则}, name={书写原则}]{writingprinciple}

% Uniformly change numbering depth
\newcommand{\numberitemswithin}[1]{%
\numberwithin{theorem}{#1}%
\numberwithin{proposition}{#1}%
\numberwithin{lemma}{#1}%
\numberwithin{corollary}{#1}%
\numberwithin{definition}{#1}%
\numberwithin{notation}{#1}%
\numberwithin{axiom}{#1}%
\numberwithin{example}{#1}%
\numberwithin{exercise}{#1}%
\numberwithin{discussion}{#1}%
\numberwithin{problem}{#1}%
\numberwithin{strategy}{#1}%
}



% Hints and solutions to exercises
% Source: LaTeX Stack Exchange
% https://tex.stackexchange.com/q/415984 (Clive Newstead)
% https://tex.stackexchange.com/a/415986 (user31729)

\newcommand{\hintref}[1]{{\color{excol}\fontfamily{bch}\bfseries  \Cref{#1}} 提示\\}

\newwrite\collectedcontentfile
\AtBeginDocument{%
  % Automatically open the file at the beginning of the document
  \immediate\openout\collectedcontentfile=\jobname.cll
}

\NewEnviron{backhint}{
  \immediate\write\collectedcontentfile{%
    \expandafter\unexpanded\expandafter{\BODY}^^J
  }%
}

\newcommand{\printhints}{%
  % Closing the file
  \immediate\closeout\collectedcontentfile% 
  \InputIfFileExists{\jobname.cll}{}{}
}

\newcommand{\hintlabel}[2]{\label{#1}
\begin{backhint}
\hintref{#1}
#2
\end{backhint}%
}

\newcommand{\hintsection}[1]{%
\begin{backhint}
\subsection*{#1}
\end{backhint}%
}